\graph[Honeycomb lattice with \SI{1}{\centi\meter} edge lengths]{
	
	\def\spacing{1cm}
	\def\width{17}
	\def\height{19.05255888} % 22 * sqrt(3)/2
	
	% Clip to page boundaries (centered on origin)
	\pgfmathsetmacro{\halfwidth}{\width/2}
	\pgfmathsetmacro{\halfheight}{\height/2}
	\clip (-\halfwidth cm, -\halfheight cm) rectangle ++ (\width cm, \height cm);
	
	% For a regular hexagon with edge length a:
	% - Circumradius R = a = 1cm
	% - Horizontal spacing between hexagon centers = sqrt(3) * a
	% - Vertical spacing between rows = 3/2 * a
	\pgfmathsetmacro{\sqrtthree}{sqrt(3)}
	\pgfmathsetmacro{\hexspacingx}{\sqrtthree * \spacing}
	\pgfmathsetmacro{\hexspacingy}{1.5 * \spacing}

    \coordinate (a) at (1.5*\spacing , \sqrtthree/2*\spacing);
    \coordinate (b) at (1.5*\spacing , -\sqrtthree/2*\spacing);
	
    \foreach \i in {-8,-7,-6,-5,-4,-3,-2,-1,0,1,2,3,4,5,6,7,8}{
        \foreach \j in {-8,-7,-6,-5,-4,-3,-2,-1,0,1,2,3,4,5,6,7,8}{
            \coordinate (C) at ($\i*(a) + \j*(b)$);
                    
                % Draw hexagon with 1cm edges centered at (xpos, ypos)
                % For a regular hexagon with edge length a, circumradius R = a = 1cm
                % Using pointy-top orientation (typical for honeycombs)
            \draw ($(C) + (0:\spacing)$)
            -- ($(C) + (60 :\spacing)$)
            -- ($(C) + (120:\spacing)$)
            -- ($(C) + (180:\spacing)$)
            -- ($(C) + (240:\spacing)$)
            -- ($(C) + (300:\spacing)$)
            -- cycle;
        }
    }
}

